%% Causal Actionable Counterfactuals --- main paper.
%% Build with `make paper` (needs a full TeX Live: pdflatex + bibtex, with the
%% booktabs, amssymb, graphicx, natbib, hyperref, algorithm, algpseudocode
%% packages). Figures are read from paper/figures/ (pdflatex prefers the .pdf
%% variants); tables are \input from paper/tables/. Both are regenerated by the
%% experiments, so this file never hard-codes a number --- every result lives in
%% a \input-ed table or a referenced figure.
\documentclass[11pt]{article}

\usepackage[margin=1in]{geometry}
\usepackage[T1]{fontenc}
\usepackage[utf8]{inputenc}
\usepackage{lmodern}
\usepackage{microtype}
\usepackage{amsmath}
\usepackage{amssymb}      % \checkmark, used by the generated tables
\usepackage{graphicx}
\usepackage{booktabs}
\usepackage{array}
\usepackage{caption}
\usepackage{subcaption}
\usepackage{algorithm}
\usepackage{algpseudocode}
\usepackage{xcolor}
\usepackage[numbers,round]{natbib}
\usepackage[hidelinks]{hyperref}

\graphicspath{{figures/}}

% A fallback so the paper still compiles if a minimal TeX lacks amssymb's glyph.
\providecommand{\checkmark}{\ensuremath{\surd}}
% Shorthands for the four audit names, so prose and headings stay consistent.
\newcommand{\Valid}{\textsc{valid}}
\newcommand{\Causal}{\textsc{causal}}
\newcommand{\Action}{\textsc{actionable}}
\newcommand{\Usable}{\textsc{usable}}
\newcommand{\cac}{\textsc{cac}}
\newcommand{\dopart}[1]{\ensuremath{\mathrm{do}(#1)}}

\title{\bfseries Causal Actionable Counterfactuals:\\[2pt]
  \large From Black-Box Rejection to Recourse That Survives Its Own Consequences}
\author{Anonymous submission}
\date{}

\begin{document}
\maketitle

\begin{abstract}
A counterfactual explanation tells a rejected applicant what to change to be
accepted. Almost every method produces one that flips the model --- and almost
none check whether the advice describes a world that could actually come about.
We separate three things a plan needs and show they are not the same: it must
\emph{flip} the black box (validity), it must still flip it once the causal
consequences of the changes are propagated (causal validity), and it must be
something a person could carry out (actionability). We call a plan \Usable{} only
when all three hold. Across four public decision datasets, every method we test
clears the validity bar, yet the \Usable{} rate spreads from as low as
\mbox{13\%} to \mbox{100\%} --- a gap that existing evaluations, which stop at
validity, never see. We introduce \emph{Causal Actionable Counterfactuals}
(\cac), which searches for an \emph{action} rather than a point and scores every
candidate on the record its structural causal model says the action would
produce, so a causally impossible plan is never even a candidate. \cac{} closes
the gap to 100\% usable recourse on all four datasets; an ablation shows the
causal term is what closes it, and it does so at a modest and honestly-reported
increase in effort. All datasets are public, all models run locally, and the
pipeline is fully reproducible from a single command.
\end{abstract}

\section{Introduction}
\label{sec:intro}

When an automated system denies someone a loan, a credit line, or a job, the
person is owed more than a score: they are owed an explanation they can
\emph{act} on. The dominant formalisation of such an explanation is the
\emph{counterfactual} of \citet{wachter2017counterfactual} --- the nearest input
$x'$ that the model would classify differently --- and a large literature now
generates these efficiently and diversely
\citep{mothilal2020dice,vanlooveren2021interpretable,verma2020counterfactual}.
Almost all of it optimises, and reports, a single property: \emph{validity}, that
$x'$ flips the model.

Validity is necessary. It is very far from sufficient. A counterfactual is a
piece of advice, and advice can flip a model while describing a world that could
never come about. We see all three of the following in the plans real methods
produce on standard data (\S\ref{sec:experiments}). A method tells an applicant
to acquire \$28{,}377 in capital gains --- valid, and useless. Another offers
additional capital \emph{losses} as a route to a higher income, exploiting a
correlation the model learned but no person can act on. A third flips the
decision by having the applicant work \emph{fewer} hours. Each is a valid
counterfactual under the standard definition. None is recourse.

These failures are not noise; they are two distinct, systematic gaps that
validity cannot see. The first is \emph{causal}: features cause one another, so
intervening on one moves others, and the record that actually results after a
person acts is not the record the method scored
\citep{karimi2021recourse,pearl2009causality}. Advice that ignores this is scored
in a world that acting on it destroys. The second is \emph{actionability}: a plan
may ask a person to change something immutable, protected, or only movable in one
direction \citep{ustun2019actionable,poyiadzi2020face}. Both gaps have been
studied in isolation; what has been missing is a single evaluation that holds a
plan to all of it at once, and a method built to pass.

\paragraph{The usable gap.} We separate three things a plan needs and treat them
as independent audits. A plan is \Valid{} if it flips the black box; \Causal{} if
it still flips the box once the structural causal model propagates the changes
into their consequences; and \Action{} if a person could actually carry it out.
We call a plan \Usable{} only when all three hold. The three are not nested: a
plan can be actionable yet causally invalid, or causally valid yet ask the
impossible. Measuring them separately is the point --- it turns ``does it flip the
model'' into ``is this advice you could follow and would it work,'' and it
reveals a gap the field has been reporting past. Figure~\ref{fig:usable_gap}
shows it: on our four datasets every method clears the validity bar, and the
\Usable{} rate nonetheless ranges from \mbox{13\%} to \mbox{100\%}.

\begin{figure}[t]
  \centering
  \includegraphics[width=0.86\linewidth]{results/results_usable_gap}
  \caption{The validity--usable gap. Faint bars: the share of rejected applicants
  a method gives a \emph{valid} plan (it flips the black box). Solid bars: the
  share it gives a \Usable{} plan (valid \emph{and} causally sound \emph{and}
  actionable). Every method clears validity; only \cac{} closes the gap on all
  four datasets. Numbers in Table~\ref{tab:main_results}.}
  \label{fig:usable_gap}
\end{figure}

\paragraph{Causal Actionable Counterfactuals.} We then close it. \cac{} searches
for an \emph{action} --- a small set of interventions $\dopart{X_S = v_S}$ --- and
scores each candidate not on the point it names but on the record its structural
causal model says the action would \emph{produce}, with every downstream feature
recomputed. Causal validity thereby becomes a property of the representation
rather than a penalty that can be traded away: \cac{} cannot propose a world its
own model disowns. Around that core it jointly minimises the effort of the
actions, their feasibility under an explicit actionability model, and the
plausibility of the resulting record, with a population search because the black
box is opaque and half the search space is a discrete subset choice.

\paragraph{Contributions.}
\begin{itemize}
  \item \textbf{A three-audit evaluation of recourse} (\S\ref{sec:desiderata})
    that separates validity, causal validity, and actionability, and scores
    methods on their conjunction, \Usable{} --- exposing a gap that
    validity-only evaluation hides (Figure~\ref{fig:usable_gap},
    Table~\ref{tab:main_results}).
  \item \textbf{The \cac{} method} (\S\ref{sec:method}): a gradient-free search
    over interventions, scored through a structural causal model, that is
    causally valid by construction and jointly optimises cost, feasibility, and
    plausibility.
  \item \textbf{A reproducible study on four public datasets}
    (\S\ref{sec:experiments}) showing the gap is real and large, that \cac{}
    closes it to 100\% usable recourse everywhere, that an ablation attributes
    the closing to the causal term, and that the cost of doing so is modest and
    honestly reported. Every dataset is public, every model runs locally, and the
    whole pipeline rebuilds from one command.
\end{itemize}

\section{Related work}
\label{sec:related}

\paragraph{Counterfactual explanations.} \citet{wachter2017counterfactual}
framed a counterfactual as the closest input that changes the decision, found by
minimising a distance-to-target objective. The idea was extended to produce
\emph{diverse} sets \citep{mothilal2020dice}, to stay on the data manifold with
prototypes \citep{vanlooveren2021interpretable} or density-weighted paths
\citep{poyiadzi2020face}, and to trade objectives explicitly with multi-objective
search \citep{dandl2020multi}. Surveys catalogue dozens of such methods
\citep{verma2020counterfactual,karimi2022survey} and benchmarks compare them
\citep{pawelczyk2021carla,klaise2021alibi}. The near-universal headline metric is
validity, sometimes with proximity and diversity; the question of whether the
named point could be \emph{reached} is largely left aside, and the assumptions
that licence reading a point as advice --- above all that its features can be set
independently of one another --- typically go unstated \citep{barocas2020hidden}.

\paragraph{Actionable recourse.} \citet{ustun2019actionable} made mutability
first-class: some features cannot be changed, others only in one direction, and
recourse must live in the actionable subspace. This is an essential correction,
and we adopt it wholesale as one of our three audits. It is, however, a constraint
on the \emph{named} change; on its own it still treats the world as a set of
independently-settable dials.

\paragraph{Causal recourse.} \citet{karimi2021recourse} argued that recourse is
about \emph{interventions}, not points: because features are causally linked, one
should optimise over actions $\dopart{X_S = v_S}$ and read off consequences
through a structural causal model \citep{pearl2009causality}, and they show a
minimal-cost intervention can differ sharply from the nearest counterfactual.
\citet{karimi2020imperfect} extend this to imperfect causal knowledge. Our method
sits squarely in this line --- it is the mechanism by which \cac{} is causally
valid by construction. Our contribution relative to it is (i) an \emph{evaluation}
that scores existing point-based methods on causal validity and actionability
together, quantifying how far short they fall, and (ii) a single search that adds
an explicit human-effort cost, an actionability guard, and a plausibility term
around the causal core, rather than causal validity alone.

\paragraph{Robustness.} A separate concern is whether recourse survives model
change \citep{upadhyay2021robust}; we touch it only as one axis of our robustness
study (\S\ref{sec:experiments}), where we re-audit \cac{} under a perturbed SCM, a
swapped black box, a tighter budget, and a different plausibility neighbourhood.

\paragraph{Position.} The pieces --- validity, actionability, causal consequence
--- exist. What has been missing, and what we supply, is (a) an evaluation that
insists on \emph{all three at once} and reports their conjunction, and (b) a
method that treats causal validity as a property of the search space rather than a
soft penalty, so it cannot be traded away for a lower distance.

\section{Problem setup}
\label{sec:problem}

\paragraph{The black box.} A trained classifier $f:\mathcal{X}\to[0,1]$ maps an
instance $x=(x_1,\dots,x_d)$ to a score, the estimated probability of the
favourable outcome (loan granted, credit line extended, low diabetes risk). A
threshold $\tau$ turns the score into a decision
$\hat y(x)=\mathbb{1}[f(x)\ge\tau]$. We treat $f$ as \emph{sacred}: the only
things any method --- ours included --- may ask of it are $f(x)$ and the induced
$\hat y(x)$. No gradients, no trees, no feature importances. This is both a
faithfulness requirement (real deployed models are often unavailable for
inspection) and a discipline that keeps the comparison honest, since the
baselines that can peek would otherwise be solving a different problem.

Fix an applicant $x^{0}$ who was rejected, $f(x^{0})<\tau$. Recourse is advice
that moves them to the favourable side.

\paragraph{Two things called ``a counterfactual.''} The literature and this paper
use the word for two different objects, and the gap between them is the whole
story.
\begin{itemize}
  \item A \emph{counterfactual point} $x'$ is an instance with
    $f(x')\ge\tau$. This is what distance-minimising methods return: a
    destination.
  \item An \emph{action} $a=\{(j,v_j):j\in S\}$, $S\subseteq\{1,\dots,d\}$,
    $|S|\le K$, is a set of interventions $\dopart{X_S=v_S}$ the applicant
    performs on features they control. This is a route.
\end{itemize}
A point is not a route. Setting $x'_j$ for the features a method chose to change
does not leave the other features where the method left them, because features
cause one another.

\paragraph{The world an action produces.} We equip each dataset with a structural
causal model $\mathcal{M}$ (\S\ref{sec:scm}): a directed acyclic graph over the
features and, for each feature, a structural equation
$X_j := g_j\big(\mathrm{pa}(j),U_j\big)$ from its parents and an exogenous term.
Given the applicant $x^{0}$ and an action $a$, the record that actually results is
obtained by the standard abduction--action--prediction steps
\citep{pearl2009causality}: infer the exogenous terms consistent with $x^{0}$,
apply $\dopart{X_S=v_S}$, and recompute every descendant of $S$ through its
structural equation. Write the result
\begin{equation}
  x^{a} \;=\; \mathcal{M}\big(x^{0};\,\dopart{X_S=v_S}\big).
  \label{eq:propagate}
\end{equation}
Features that are not descendants of $S$ keep their values; descendants move to
whatever the equations imply. It is $x^{a}$, not the named point, that the world
presents to the black box after the applicant acts.

\paragraph{The mismatch, stated plainly.} A point-based method reports success
when $f(x')\ge\tau$. But the applicant, acting on the changes in $x'$, arrives at
$x^{a}$, and the decision they actually receive is $\hat y(x^{a})$. When
$f(x')\ge\tau$ but $f(x^{a})<\tau$, the advice flipped the model on paper and
failed the person in fact. Section~\ref{sec:desiderata} turns this mismatch, and
the separate question of whether $a$ is something a person could even do, into
measurable audits; Section~\ref{sec:method} builds a method that scores $x^{a}$
directly and so cannot fall into the gap.

\section{What makes recourse usable}
\label{sec:desiderata}

We hold every plan --- whether a point $x'$ from a baseline or an action $a$ from
\cac{} --- to three audits. Each is a yes/no property; a plan is \Usable{} exactly
when it passes all three. To audit a baseline that returns a point, we read the
point as the action ``set the features it changed to the values it named'' and
propagate it through Eq.~\eqref{eq:propagate} like any other action, so every
method is judged on the world its advice produces.

\paragraph{1. Validity ($\Valid$).} The plan flips the black box \emph{as named}:
\begin{equation}
  \Valid(x') \;=\; \mathbb{1}\big[f(x')\ge\tau\big].
\end{equation}
This is the field's usual metric, and the necessary floor. For \cac{} the named
point is $x^{a}$ itself, so \cac{} passes validity iff it passes the next audit.

\paragraph{2. Causal validity ($\Causal$).} The plan still flips the box once its
changes are propagated into their consequences:
\begin{equation}
  \Causal(a) \;=\; \mathbb{1}\big[f(x^{a})\ge\tau\big],
  \qquad x^{a}=\mathcal{M}\big(x^{0};\dopart{X_S=v_S}\big).
\end{equation}
A point-based plan can pass \Valid{} and fail \Causal{}: it named a destination
that the causal dynamics never deliver. This is the audit the field skips.

\paragraph{3. Actionability ($\Action$).} A person could carry the plan out. Using
the per-dataset actionability model of \S\ref{sec:cost}, every acted-on feature
$j\in S$ must be mutable, must move in an admissible direction (income may rise,
age may not be lowered, a diagnosis is fixed), and must stay within its plausible
bounds; and the plan must touch no protected attribute:
\begin{equation}
  \Action(a) \;=\; \prod_{j\in S}\;
    \mathbb{1}\big[\text{$j$ mutable}\big]\,
    \mathbb{1}\big[\text{direction of $v_j$ admissible}\big]\,
    \mathbb{1}\big[v_j\in[\ell_j,u_j]\big]\,
    \mathbb{1}\big[\text{$j$ not protected}\big].
\end{equation}
A plan can be causally valid yet fail here (it works, but asks the impossible),
or actionable yet causally invalid (a person could do it, but it would not
land).

\paragraph{The conjunction.} A plan is usable iff it survives all three:
\begin{equation}
  \Usable(a) \;=\; \Valid(x^{a})\;\wedge\;\Causal(a)\;\wedge\;\Action(a).
  \label{eq:usable}
\end{equation}
The three are logically independent --- each of the three ``valid but not usable''
plans quoted in \S\ref{sec:intro} fails a different one --- so their conjunction
is a genuinely stricter bar than any single audit, and it is the bar the paper
reports.

\paragraph{Crediting a method.} A method may offer several ranked suggestions per
applicant (we take up to five). We credit it with usable recourse for an applicant
if \emph{any} of its suggestions is usable, mirroring how a person would use the
tool: they need one plan they can follow, not five. A method's \Usable{} rate in
Table~\ref{tab:main_results} is the fraction of rejected applicants for whom it
clears this bar. Per-suggestion audit records are retained for the appendix and
the ablation.

\paragraph{Two quality measures (not gates).} Among usable plans, two more
properties decide which advice is \emph{better}, and we report them without
letting them gate usability.
\emph{Effort} (cost, \S\ref{sec:cost}) is the normalised human difficulty of the
actions; lower is kinder.
\emph{Plausibility} is the distance from $x^{a}$ to the favourable-class data,
measured as the mean distance to its $k$ nearest favourable neighbours and
normalised by the dataset's median such distance; lower means the resulting record
looks like a real accepted applicant rather than an adversarial artefact. These
let us ask not just whether a method gives usable recourse but whether the recourse
is cheap and realistic --- the subject of the cost-of-validity analysis in
\S\ref{sec:experiments}.

\section{Datasets and black boxes}
\label{sec:datasets}

We study four public tabular decision problems, chosen because each is a setting
where a rejected person would genuinely want recourse and where the features carry
real causal structure.

\begin{itemize}
  \item \textbf{Adult} \citep{kohavi1996adult,dua2019uci}: predict whether income
    exceeds \$50K. We cap the data at $\text{MAX\_ROWS}=20{,}000$ rows for
    tractability and hold out 20\% ($n=4000$) for testing. The favourable outcome
    is high income; \texttt{gender} is protected and \texttt{age} immutable.
  \item \textbf{German Credit} \citep{dua2019uci}: 1{,}000 applicants, test
    $n=200$; predict good versus bad credit risk. The favourable outcome is good
    credit; \texttt{age} and \texttt{sex} are protected.
  \item \textbf{HELOC} \citep{fico2018heloc}: home-equity line-of-credit
    applications, test $n=1975$; predict whether repayment risk is acceptable. The
    features are dense, monotone credit-bureau aggregates, which makes the causal
    structure both important and hard to get exactly right (\S\ref{sec:scm}).
  \item \textbf{Pima diabetes} \citep{smith1988diabetes}: 768 records, test
    $n=154$; predict diabetes onset. Because recourse should move a person toward
    the \emph{healthy} outcome, we invert the label so the favourable class is
    \emph{no diabetes} ($\texttt{no\_diabetes}=1-\texttt{Outcome}$).
    \texttt{Glucose}, \texttt{BMI}, and \texttt{BloodPressure} are actionable;
    \texttt{Age}, \texttt{Pregnancies}, and \texttt{DiabetesPedigreeFunction} are
    not. We include it as the one \emph{health} setting alongside the three credit
    and income problems, to check that the usable gap is not a peculiarity of
    financial data.
\end{itemize}

\paragraph{The black box.} For every dataset the primary black box is an XGBoost
classifier \citep{chen2016xgboost} and the secondary is a multilayer perceptron,
each trained on the 80\% training split with \texttt{random\_state}$=42$ and a
decision threshold of $0.5$. Table~\ref{tab:blackbox_performance} reports held-out
performance; the primary XGBoost model clears its pre-registered AUC target on all
four datasets (Adult 0.918, Diabetes 0.849, HELOC 0.788, German 0.785). Neither
model is inspected by any recourse method: the whole study treats $f$ through
\texttt{predict} and \texttt{predict\_proba} alone (\S\ref{sec:problem}).

\paragraph{The query set.} Recourse is only meaningful for someone who was
rejected, so for each dataset the \emph{query set} is the test records the primary
black box rejects, capped at $\text{N\_QUERY\_INSTANCES}=100$: Adult and HELOC
hit the cap at 100, German has 51, diabetes 43. Every method is asked for up to
$k=\text{N\_COUNTERFACTUALS}=5$ suggestions per query record, and every table in
\S\ref{sec:experiments} is computed on this shared set.

\input{tables/blackbox_performance}

\section{Structural causal models}
\label{sec:scm}

Causal validity is only as meaningful as the causal model behind it. For each
dataset we hand-craft a directed acyclic graph over the features from domain
knowledge --- education raises income, age drives credit history, glucose and BMI
move together --- and fit a \emph{functional} causal model: every non-root feature
is regressed on its parents and its fitted residual is retained as the exogenous
term, so the model can run the abduction--action--prediction of
Eq.~\eqref{eq:propagate} on any record. The graphs and fitted equations live in
\texttt{models/fcm\_*.joblib}; Figure~\ref{fig:scm_adult} draws the Adult graph.

\begin{figure}[t]
  \centering
  \includegraphics[width=0.72\linewidth]{scm_adult}
  \caption{The Adult structural causal model: 10 nodes, 25 edges, 7 of the
  features actionable. Arrows are asserted causal influences; the target node is
  reachable from the actionable set, which is what recourse requires.}
  \label{fig:scm_adult}
\end{figure}

\paragraph{These graphs are approximations, and we test them.} A hand-drawn graph
is a modelling choice, not ground truth, and honesty about how far the data back
it is part of the contribution. We run two tests
(Table~\ref{tab:scm_validation}). An \emph{edge test} asks whether each asserted
arrow is still detectable given the child's other parents; an
\emph{independence test} asks whether each arrow the graph \emph{omits} is
consistent with the data. Both use Benjamini--Hochberg-corrected $q\le0.05$ and a
minimum partial $R^2$ of $0.01$ on 5{,}000 training rows.

The picture is mixed and we report it in full. All four graphs are acyclic with
the target reachable from the actionable set --- the structural precondition
\cac{} needs. The diabetes graph is the cleanest (16 of 17 edges supported, all 27
implied independencies upheld); Adult is solid (17 of 25 edges, one outright
unsupported). HELOC and especially German are weaker: German supports only 18 of
37 asserted edges and upholds 88 of 138 independencies, because credit features
are densely inter-correlated and a parsimonious graph inevitably omits real
dependences among them. We keep these graphs rather than densifying them to chase
the tests, and we treat the gap as a limitation --- one the robustness study meets
head-on by re-scoring \cac{} under deliberately perturbed and even fully
disconnected graphs (\S\ref{sec:experiments}). The per-edge results are in
Appendix~\ref{app:scm}. The Adult graph as used is a 10-node, 25-edge superset of
the smaller graph sketched in the specification; this, and the imperfect German
and HELOC graphs, are recorded among the README deviations.

\section{The cost of an action}
\label{sec:cost}

Actionability and effort are two questions about the same plan --- \emph{may} a
person take this action, and \emph{how hard} is it --- and \S\ref{sec:desiderata}
uses both: the first gates $\Action$, the second is the \emph{cost} quality
measure. We answer them with one hand-crafted model per dataset
(Table~\ref{tab:cost-model-adult} shows Adult; the rest are in
Appendix~\ref{app:cost}).

\paragraph{What is permitted.} Each feature carries an admissibility rule: whether
it is mutable at all, the only direction it may move (income may rise, age may not
fall, a past diagnosis is fixed), and the bounds it must stay within. Protected
attributes (\texttt{gender}, \texttt{sex}, \texttt{age}) and structurally fixed
ones are marked immutable. An action that touches an immutable or protected
feature, or pushes one the wrong way, is inadmissible --- it fails $\Action$ and,
equivalently, is priced at \texttt{INFINITE\_COST} in the search.

\paragraph{How hard it is.} For each mutable feature we specify an effort function
of the amount moved --- \texttt{constant} for a lump change, a \texttt{ladder} for
ordinal steps (each education level costs more than the last), a \texttt{switch}
for unordered categories, or a \texttt{piecewise} rate that rises as the move
grows (the first ten hours a week are cheap, the next twenty dear). Each is scaled
by a weight read from the SCM's qualitative effort label for that feature
(\texttt{medium}, \texttt{high}, \texttt{very\_high}, \dots), so the graph's sense
of what is hard and the numeric cost stay consistent. The cost of a multi-feature
action is the sum over its interventions.

\subsection{The budget}
\label{sec:budget}
Raw effort is not comparable across datasets, so we normalise. The
\emph{budget} for a dataset is the summed cost of its
$\text{MAX\_ACTIONS}=3$ most expensive admissible single actions --- the effort of
the hardest plan the rules still allow --- and every plan's cost is reported as a
fraction of it, so any permitted plan lands in $[0,1]$. For Adult the budget is
$11.22$. This normalisation is what lets Table~\ref{tab:main_results} put effort
on one scale across four very different problems; it also sets the search ceiling
and the audit's ``too many actions'' rule, which is why the budget-sensitivity
study (\S\ref{sec:experiments}) varies all three together. Two ratio views appear
in the results --- cost against the cheapest \emph{oracle} flip on the same record
(Table~\ref{tab:main_results}) and against \cac{}'s own full objective
(Table~\ref{tab:ablation}); the README notes both definitions.

\section{Causal Actionable Counterfactuals}
\label{sec:method}

\cac{} makes causal validity structural rather than optional. It never names a
destination point; it searches for an \emph{action}, and it scores that action on
the record its causal model says the action produces. The audited object and the
optimised object are therefore the same record $x^{a}$ of
Eq.~\eqref{eq:propagate}, so a plan that would not survive its own consequences is
not a low-scoring candidate --- it is not the plan \cac{} is scoring at all.

\subsection{Objective}
A candidate action $a$ acts on a feature subset $S$ with $|S|\le K$
($K=\text{MAX\_ACTIONS}=3$) and is scored through the resulting record
$x^{a}=\mathcal{M}(x^{0};\dopart{X_S=v_S})$ by
\begin{equation}
  J(a) \;=\; \lambda_{\text{v}}\,\underbrace{\big[\tau+\delta-f(x^{a})\big]_{+}}_{\text{validity}}
        \;+\; \lambda_{\text{c}}\,\underbrace{c(a)}_{\text{cost}}
        \;+\; \lambda_{\text{f}}\,\underbrace{\phi(a)}_{\text{feasibility}}
        \;+\; \lambda_{\text{p}}\,\underbrace{\rho(x^{a})}_{\text{plausibility}},
  \label{eq:objective}
\end{equation}
minimised over $a$, with $[\,\cdot\,]_{+}=\max(0,\cdot)$ and weights
$(\lambda_{\text{v}},\lambda_{\text{c}},\lambda_{\text{f}},\lambda_{\text{p}})
=(1.0,\,0.6,\,0.3,\,0.2)$.
\begin{itemize}
  \item \textbf{Validity} is a hinge on the target score that vanishes once
    $f(x^{a})$ clears the threshold $\tau$ by a margin
    $\delta=\text{CAC\_TARGET\_MARGIN}=0.05$. Because it is evaluated on the
    propagated record, driving this term to zero \emph{is} causal validity ---
    there is no separate ``flip the named point'' to fall short of.
  \item \textbf{Cost} $c(a)\in[0,1]$ is the normalised effort of the actions
    (\S\ref{sec:cost}), priced on what the person does, not on the consequences
    the model adds.
  \item \textbf{Feasibility} $\phi(a)$ adds the \texttt{INFINITE\_COST}
    admissibility guard (an inadmissible action is unbounded here, so it can never
    win) and a penalty for disturbing sensitive attributes downstream.
  \item \textbf{Plausibility} $\rho(x^{a})$ is the mean distance from $x^{a}$ to
    its $k{=}5$ nearest favourable-class neighbours, normalised by the dataset
    median, so the produced record resembles a real accepted applicant.
\end{itemize}

\subsection{Search}
The objective is not differentiable and cannot be made so: the black box is opaque
by construction (\S\ref{sec:problem}), the choice of which features to act on is a
discrete subset selection, and the causal model is piecewise-constant wherever a
feature is categorical. \cac{} therefore uses a gradient-free elitist population
search (Algorithm~\ref{alg:cac}), in the elitist tradition of
\citet{deb2002nsga}. A candidate is a mask marking the acted-on
features together with their target values; the population is evaluated in one
batch through the causal model, the better half is retained, and the rest are
regenerated by crossover, Gaussian jitter on the values, and a set mutation that
adds or drops a feature from the acting set. Every child is \emph{legalised} ---
clipped to admissible bounds and grids, stripped of immutable features, and capped
at $K$ actions --- so the search only ever scores plans a person could take. The
$k$ best distinct survivors are returned as the applicant's ranked suggestions.

\begin{algorithm}[t]
\caption{\cac{} action search for one rejected applicant.}
\label{alg:cac}
\begin{algorithmic}[1]
\Require black box $f$, causal model $\mathcal{M}$, applicant $x^{0}$, target class,
  cost model, favourable data
\State $n \gets \max(\text{CAC\_POPULATION},\,\text{MIN\_POPULATION}) = 60$;\quad
  $m \gets \max(2,\lceil n\cdot\text{ELITE\_FRAC}\rceil)=30$
\State initialise $n$ legal candidates: random acting sets of size $\le K$ with
  values drawn from each feature's admissible grid
\For{$g = 1$ \textbf{to} $\text{CAC\_MAX\_ITERS}=200$}
  \For{each candidate $a$}\Comment{one batched pass}
    \State $x^{a} \gets \mathcal{M}(x^{0};\dopart{X_S=v_S})$
      \Comment{abduction--action--prediction}
    \State score $J(a)$ by Eq.~\eqref{eq:objective}
  \EndFor
  \State $E \gets$ the $m$ candidates of lowest $J$\Comment{elitist retention}
  \State children $\gets$ from $E$: uniform crossover, then value jitter
    $\mathcal{N}(0,\text{CAC\_LEARNING\_RATE}{=}0.1)$, then set mutation
    (prob.\ $0.35$)
  \State legalise every child (bounds, grids, mutability, $|S|\le K$)
  \State population $\gets E \cup \text{children}$
\EndFor
\State \Return the $k$ distinct legal candidates of lowest $J$
\end{algorithmic}
\end{algorithm}

\subsection{A note on reconstruction}
The specification we built from was truncated inside the subsection that gave
\cac{}'s pseudocode, so the objective of Eq.~\eqref{eq:objective} and the search of
Algorithm~\ref{alg:cac} are a faithful reconstruction consistent with the parts of
the specification that survived --- the four named terms, their weights, the
action budget, and the intervention-through-SCM scoring. This is stated plainly
among the deviations in the README so no reader mistakes the reconstruction for a
verbatim reproduction. What matters for the claims that follow is the property the
reconstruction preserves exactly: \cac{} scores interventions through the causal
model, so it cannot propose a world its own model disowns.

\section{Experiments}
\label{sec:experiments}

We ask four questions. Is the validity--usable gap real, and how large
(\S\ref{sec:exp:gap})? \emph{Why} do valid plans fail --- causally, or on
actionability (\S\ref{sec:exp:why})? Is the closing of the gap attributable to the
causal term specifically (\S\ref{sec:exp:ablation})? And what does \cac{} pay for
it, and does the advantage survive changes to the setup
(\S\ref{sec:exp:cost}--\ref{sec:exp:robust})?

\paragraph{Baselines.} Against \cac{} we run four references, all querying the same
opaque box. \textbf{Random} samples admissible perturbations and keeps those that
flip. \textbf{Wachter} \citep{wachter2017counterfactual} minimises distance to a
flipping point. \textbf{DiCE} \citep{mothilal2020dice} returns diverse flipping
points. \textbf{Oracle} returns the lowest-cost \emph{admissible} action that
flips the box, found by direct search; it defines the \emph{Cost/Oracle}
denominator and serves as a strong actionable-but-not-causal reference --- it is
the cost floor a method would hit if causal validity were free. All methods are
given $k=5$ suggestions per applicant and audited by
Eq.~\eqref{eq:usable}.

\subsection{The gap is real, and largest where structure is densest}
\label{sec:exp:gap}
Table~\ref{tab:main_results} is the headline. Every method clears validity almost
everywhere --- the column is a wall of $1.00$s, with only Wachter and Random
dipping on HELOC --- and the \Usable{} column then fans out. The spread is
dataset-dependent, and honestly so. On \textbf{Adult} the gap is small: the box's
decision leans on \texttt{hours\_per\_week} and \texttt{education}, features that
are cleanly actionable and causally shallow, so every method lands near
$1.00$ usable and \cac{}'s discipline buys little. The gap opens exactly where the
features are densely causal. On \textbf{HELOC}, DiCE is $1.00$ valid but $0.50$
usable and Wachter $0.85$ valid but $0.13$ usable; even the cost-optimal Oracle
reaches only $0.66$. On \textbf{German}, Wachter falls from $1.00$ valid to $0.45$
usable. \cac{} is $1.00$ usable on all four. Figure~\ref{fig:usable_gap} draws the
faint-to-solid drop; Figure~\ref{fig:funnel} traces it audit by audit.

\begin{figure}[t]
  \centering
  \includegraphics[width=0.86\linewidth]{results/results_recourse_funnel}
  \caption{The recourse funnel. Starting from every rejected applicant, each stage
  keeps only the plans that survive one more audit --- valid, then also causal,
  then also actionable. The baselines leak at the causal and actionability stages;
  \cac{} does not, because it is built to pass them.}
  \label{fig:funnel}
\end{figure}

\subsection{Why valid plans fail}
\label{sec:exp:why}
The two middle audits explain the fan-out.
Table~\ref{tab:faithfulness} isolates \emph{causal} failure over the plans that
flip the box: it reads each proposal as the intervention it implies, propagates it,
and re-queries $f$. On HELOC only $27\%$ of Wachter's valid plans and $36\%$ of
DiCE's are causally valid; their flip survives propagation just over half the
time, and the mean loss in target probability is \emph{positive} ($+0.111$ for
Wachter, $+0.090$ for DiCE), meaning the consequences of the advice actively push
the applicant back across the line. On German the same story holds ($42\%$, $47\%$
causal valid). \cac{} sits at $1.00$ with a zero propagation gap by construction:
the plan it proposed and the record it was scored on are the same object.

Table~\ref{tab:feasibility} isolates \emph{actionability}, on what each method
\emph{asks for}. It is where Wachter truly collapses: on HELOC only $8\%$ of its
valid plans are actionable --- $84\%$ move a feature against its only admissible
direction and it spends $5.6$ actions on average, far past the three-action
budget. DiCE is far better ($86\%$ actionable on HELOC) but still trips the
immutability rule on a tenth of plans. \cac{} and Oracle are $1.00$ actionable
everywhere by construction. The two tables together show the failures are
genuinely distinct: DiCE's problem on HELOC is mostly causal, Wachter's is both
causal and gross inactionability, and no single-audit evaluation would have caught
either.

\subsection{The causal term is what closes the gap}
\label{sec:exp:ablation}
Table~\ref{tab:ablation} and Figure~\ref{fig:ablation} remove one objective term
at a time and re-audit. The result is clean: dropping the cost, feasibility,
validity, or plausibility term leaves \cac{} at $1.00$ usable on every dataset ---
each term moves its own axis (removing cost roughly doubles effort on German,
$0.009\!\to\!0.094$; removing plausibility more than doubles the plausibility
distance on Adult, $2.40\!\to\!5.04$) without breaking usability. Only removing the
\emph{causal treatment} --- scoring the identical population as a point rather than
as an intervention, the sole way to strip causal validity from a method that
represents it structurally --- breaks it: usable falls to $0.68$ on HELOC, $0.85$
on Adult, $0.94$ on German, the propagation gap jumps from zero, and the flip
survives only $60$--$79\%$ of the time. The gap is the causal term's doing.

\begin{figure}[t]
  \centering
  \includegraphics[width=0.80\linewidth]{results/results_ablation}
  \caption{Objective ablation (Table~\ref{tab:ablation}). Removing any of the
  cost, feasibility, validity, or plausibility terms keeps \cac{} fully usable;
  removing the causal (intervention-through-SCM) treatment is the only change that
  reopens the gap.}
  \label{fig:ablation}
\end{figure}

\subsection{What causal validity costs}
\label{sec:exp:cost}
\cac{} is not free, and we report the price without flattering it
(Figure~\ref{fig:cost}). On Adult, \cac{}'s effort of $0.181$ is the
\emph{highest} of any method --- above DiCE's $0.075$ and Random's $0.044$ ---
because insisting a plan survive its consequences rules out the cheap flips the
others take; in exchange its produced records are the most plausible ($2.40$
versus DiCE's $3.95$, lower is better). Elsewhere the premium is modest or absent:
\cac{} is $0.031$ on HELOC and $0.009$ on German while being the only method that
is fully usable there. Against the cost-optimal Oracle flip on the same record,
\cac{}'s median effort ratio runs from $1.24\times$ (diabetes) to $5.87\times$
(German). That multiplier is the cost of causal validity and plausibility over a
bare flip --- and the Oracle rows are the rejoinder to any suspicion that one
could simply minimise cost: Oracle is $1.00$ actionable and cost-optimal, yet only
$0.66$ usable on HELOC and $0.87$ on Adult, because it too scores the point it
names.

\begin{figure}[t]
  \centering
  \includegraphics[width=0.80\linewidth]{results/results_cost_of_validity}
  \caption{The cost of validity. \cac{} buys usable, causally-sound recourse at a
  modest and dataset-dependent increase in effort, and returns the most plausible
  post-recourse records. On Adult it is the most expensive method; on HELOC and
  German it is among the cheapest while being the only fully-usable one.}
  \label{fig:cost}
\end{figure}

\subsection{One applicant, five methods}
Table~\ref{tab:case-study} grounds the aggregate in a single case:
Adult test record~9, chosen as the easy case where \emph{every} method is usable
($86$ of $100$ applicants fall here). All five tell the applicant to work more
hours; they differ only in how many (\cac{} $+10$, Oracle $+5$, DiCE $+19$,
Wachter and Random $+3$) and thus in effort. This is the honest picture on Adult:
when the decisive feature is cleanly actionable and causally shallow, everyone
agrees, and \cac{}'s value is insurance rather than rescue. The appendix
(Table~\ref{tab:qualitative-cases}) prints the harder curated cases --- including
HELOC and German applicants where only \cac{} produces a usable plan --- action by
action.

\subsection{The advantage survives changes to the setup}
\label{sec:exp:robust}
Four robustness studies re-audit \cac{} under a moved goalpost.
\textbf{Budget} (Table~\ref{tab:robustness_budget}): with a one-action budget
\cac{} still reaches $0.95$ usable on Adult and $0.85$ on HELOC, and $1.00$ once
two actions are allowed --- three is comfortable headroom, not a crutch.
\textbf{Routes asked for} (Table~\ref{tab:robustness_k}): \cac{} is $1.00$ usable
at $k=1$ on every dataset, so its lead does not come from being handed more
attempts; DiCE needs $k=5$ to catch up on the easy datasets and still caps at
$0.55$ usable on HELOC even at $k=10$.
\textbf{Model transfer} (Table~\ref{tab:robustness_model}): \cac{}'s plans,
re-scored against the secondary MLP and against bootstrap refits of the deployed
box, stay usable for $0.85$--$1.00$ of applicants, while the baselines erode
sharply --- Wachter's HELOC recourse drops to $0.00$ usable under every audit box.
\textbf{Graph misspecification} (Table~\ref{tab:robustness_scm}): when \cac{}
searches with a deliberately corrupted graph but is scored against the true one,
single-edge errors (missing, spurious, or reversed) leave causal validity at
$0.98$--$1.00$; only erasing the structure entirely reopens the gap ($0.95$ on
Adult, $0.65$ on HELOC), which is exactly the \emph{-- SCM} ablation row. \cac{}
degrades gracefully with the quality of the causal knowledge it is given, rather
than depending on the graph being perfect --- the appropriate behaviour given the
imperfect German and HELOC graphs of \S\ref{sec:scm}.

\section{Discussion and conclusion}
\label{sec:conclusion}

A counterfactual explanation is advice, and advice is only worth giving if it can
be followed and would work. The field has largely measured whether advice flips
the model. We argued that this conflates three separate things --- flipping the
model, surviving the causal consequences of the change, and being something a
person could actually do --- and that only their conjunction, \Usable{}, is what a
rejected applicant needs. Measuring the three separately exposes a gap that
validity-only evaluation cannot see: across four public datasets every method
clears validity, yet usable recourse ranges from $13\%$ to $100\%$, and the gap
opens precisely where the features are densely causal. \cac{} closes it by
treating causal validity as structural --- it searches for interventions and
scores them through a structural causal model, so it can only ever propose a world
its own model endorses --- and an ablation shows the causal term is specifically
what closes it, at a modest and honestly-reported cost in effort.

\paragraph{Limitations.} We are candid about three. \emph{The causal models are
hand-crafted and imperfect}: the diabetes graph is well supported by the data, but
the German and HELOC graphs leave many real dependences unmodelled
(\S\ref{sec:scm}), and \cac{}'s guarantees are relative to the graph it is given
--- mitigated, not erased, by the graceful degradation we measured under graph
perturbation (\S\ref{sec:exp:robust}), and by the model-transfer study showing the
advantage does not depend on scoring against the same box \cac{} searched.
\emph{The human-feasibility audit is a model, not people}: we encode mutability
and effort from domain knowledge and ship a user-study protocol (Appendix and
\texttt{user\_study/}) to test those assumptions with real respondents, but we do
not run it here. \emph{Recourse is single-shot}: we do not model the applicant
taking one action, being re-scored, and acting again.

\paragraph{Reproducibility.} Every dataset is public and fetched from Kaggle or
\texttt{ucimlrepo}; every model trains locally; no result depends on a paid
service or a network call at evaluation time. The pipeline is deterministic
(\texttt{random\_state}$=42$) and rebuilds end to end --- data, models, SCMs, cost
models, all methods, every table and figure in this paper, and the demo --- from a
single \texttt{make all}, with \texttt{make paper} producing this PDF.

\paragraph{Outlook.} The ingredients of good recourse --- validity, causal
consequence, actionability --- were already in the literature separately. The step
that mattered was refusing to trade them off: auditing all three at once, and
building a method whose search space cannot contain a causally impossible plan.
The result is recourse that survives its own consequences. We hope the three-audit
evaluation, more than the method, is what the field takes forward: it is cheap to
adopt, it applies to any recourse method, and it turns ``does it flip the box'' back
into the question that matters --- ``could this person follow your advice, and
would it work.''


\bibliographystyle{plainnat}
\bibliography{references}

\appendix
\section{Every plan, priced action by action}
\label{app:plans}
Table~\ref{tab:main_results} scores each method on whether \emph{any} of its
suggestions is usable; the case study in \S\ref{sec:experiments} shows one
applicant. For completeness this appendix prints the full curated set: every
method's leading plan for each qualitative case, with the actions it asks for
and the audit each plan passes or fails.
\input{tables/qualitative_cases}

\section{Per-dataset structural-causal-model validation}
\label{app:scm}
\S\ref{sec:scm} summarises how far the data backs each graph. The per-dataset
edge tests are collected here.
\input{tables/scm_edge_tests_adult}
\input{tables/scm_edge_tests_german}
\input{tables/scm_edge_tests_heloc}
\input{tables/scm_edge_tests_diabetes}

\section{Cost models}
\label{app:cost}
The normalised effort of \S\ref{sec:cost}, per dataset: the per-feature unit
cost, admissible direction, and whether a feature may be acted on at all.
\input{tables/cost_model_adult}
\input{tables/cost_model_german}
\input{tables/cost_model_heloc}
\input{tables/cost_model_diabetes}

\end{document}
